% =====================================================================
%  CHAPITRE 15 — TECHNIQUES DE L'ORAL : PARLER SANS TRADUIRE
% =====================================================================
\chapter{Parler sans traduire : les techniques de fluidité}

Connaître des mots et des règles ne suffit pas : il faut \textbf{produire en temps réel}.
C'est un sport, pas un savoir. Ce chapitre donne les techniques qui transforment un savoir
passif en parole fluide.

\section{Le problème réel : le délai de traduction}
Quand vous traduisez du français dans votre tête, votre cerveau exécute cinq opérations :
comprendre le français \ra choisir les mots \ra les traduire \ra conjuguer \ra prononcer.
Résultat : trois à cinq secondes de silence, une intonation française, et des calques.
L'objectif de ce chapitre est de réduire ce circuit à \textbf{une seule opération} :
penser directement en anglais.

\begin{tableau}[Les cinq niveaux de la pensée en anglais --- progressez dans l'ordre]
\begin{tabularx}{\linewidth}{@{}p{1.4cm}p{4.0cm}X@{}}
\entete{Niv.} & \entete{Exercice} & \entete{Détail} \tabularnewline
\hline
\textbf{1} & Nommer le monde & 2 minutes, trois fois par jour : nommez en anglais tout ce
que vous voyez. \en{table, phone, window, a red car, two children, a man walking}. Aucune
traduction autorisée --- si vous ne savez pas, contournez (voir p. suivante). \tabularnewline
\textbf{2} & Phrases simples sur soi & \en{I'm tired. I want tea. It's hot. I have work at
two.} --- 20 phrases par jour, à voix haute, en marchant. \tabularnewline
\textbf{3} & Narration de l'action en cours & Commentez ce que vous faites :
\en{I'm cutting the onions. Now I'm adding oil. It smells good.} \tabularnewline
\textbf{4} & Récit du passé et projets & Le soir : 3 minutes sur votre journée. Le matin :
3 minutes sur ce que vous allez faire. \tabularnewline
\textbf{5} & Pensée abstraite et débat & Pensez directement en anglais quand vous réfléchissez
à un problème, préparez un argument, imaginez une réponse à une objection. \tabularnewline
\end{tabularx}
\end{tableau}

\begin{memo}[Le déclencheur visuel]
Choisissez un objet que vous voyez cent fois par jour --- votre téléphone, une porte, votre
montre. Chaque fois que vous le voyez, dites une phrase en anglais (même dans votre tête).
C'est un « rappel » automatique qui installe l'anglais dans votre journée sans effort
d'organisation.
\end{memo}

\section{Le contournement : parler sans connaître le mot}
La compétence la plus utile d'un locuteur intermédiaire n'est pas le vocabulaire : c'est
\textbf{la capacité à expliquer un mot manquant}. C'est ce que font les natifs eux-mêmes
--- ils disent \en{thingy}, \en{whatchamacallit}, \en{you know, the thing\ldots}

\begin{tableau}[Les sept stratégies de contournement]
\begin{tabularx}{\linewidth}{@{}p{3.6cm}p{4.2cm}X@{}}
\entete{Stratégie} & \entete{Formule} & \entete{Exemple (mot manquant : \en{corkscrew})} \tabularnewline
\hline
Définition fonctionnelle & \en{It's a thing you use to\ldots} & \en{It's a thing you use to open a wine bottle.} \tabularnewline
Description & \en{It's like a small metal tool with a spiral.} & \en{It's made of metal, it's small, it has a spiral.} \tabularnewline
Catégorie + précision & \en{It's a kind of\ldots} & \en{It's a kind of kitchen tool.} \tabularnewline
Synonyme simple & Remplacez le mot difficile par un mot facile & \en{a bottle opener} (pour du vin) \tabularnewline
Exemple d'usage & \en{You need it when you\ldots} & \en{You need it when you buy a bottle of wine.} \tabularnewline
Geste et mimique & Mimez, montrez, dessinez & (tourner le poignet) \tabularnewline
Demande d'aide & \en{How do you say\ldots?} / \en{What's the word for\ldots?} & \en{What's the word for the thing you open wine with?} \tabularnewline
\end{tabularx}
\end{tableau}

\begin{astuce}[Entraînement « mot interdit » --- 5 min/jour]
Choisissez 5 mots courants (\en{chair, water, phone, car, bread}) et \textbf{interdisez-vous
de les prononcer} pendant 3 minutes de description. Expliquez chacun sans le nommer :
\en{``It's a thing with four legs, you sit on it at the table.''} C'est exactement
l'exercice que font les interprètes. En trois semaines, vous ne serez plus jamais bloqué.
\end{astuce}

\begin{tableau}[12 phrases qui vous donnent du temps pour réfléchir]
\begin{tabularx}{\linewidth}{@{}p{6.4cm}X@{}}
\entete{Phrase} & \entete{Effet} \tabularnewline
\hline
\en{That's a good question --- let me think.} & Gagne 3 secondes, flatte l'interlocuteur. \tabularnewline
\en{Well, it depends\ldots} & Gagne du temps en ouvrant une nuance. \tabularnewline
\en{How can I put this\ldots} & Signale que vous cherchez la formulation juste. \tabularnewline
\en{What I mean is\ldots} & Reformule après une hésitation. \tabularnewline
\en{Let me see\ldots} & Neutre, très courant. \tabularnewline
\en{To be honest, I'm not sure, but I'd say\ldots} & Répond sans savoir --- bien mieux que le silence. \tabularnewline
\en{Sorry, I've lost my thread. Where was I?} & Reprend après un trou de mémoire. \tabularnewline
\en{It's on the tip of my tongue.} & Le mot est presque là. \tabularnewline
\en{Give me a second, I'll find the word.} & Annonce la recherche. \tabularnewline
\en{In French we say\ldots\ --- how do you say it in English?} & Transforme le trou en échange. \tabularnewline
\en{Anyway, the main point is\ldots} & Reprend le fil. \tabularnewline
\en{Sorry, could I start again?} & Recommencer proprement. \tabularnewline
\end{tabularx}
\end{tableau}

\section{Les \emph{fillers} naturels : remplir sans bégayer}
Les natifs ne parlent pas sans pause : ils remplissent leurs pauses avec des mots-outils.
Les connaître vous évite le « euh » français (qui sonne étranger) et vous donne l'air naturel.

\begin{tableau}[25 \emph{fillers} et connecteurs oraux]
\begin{tabularx}{\linewidth}{@{}p{3.4cm}p{3.0cm}X@{}}
\entete{\en{Filler}} & \entete{Équivalent} & \entete{Emploi et exemple} \tabularnewline
\hline
\en{well} & « eh bien » & Pour commencer une réponse : \en{Well, I think it's fine.} \tabularnewline
\en{you know} & « tu vois » & Recherche d'accord : \en{It's hard, you know?} \tabularnewline
\en{I mean} & « je veux dire » & Reformulation : \en{I mean, it's not that bad.} \tabularnewline
\en{kind of / sorta} & « un peu, genre » & Adoucissement : \en{I'm kind of tired.} \tabularnewline
\en{actually} & « en fait » & Correction d'attente : \en{Actually, I've never been there.} \tabularnewline
\en{basically} & « en gros » & Résumé : \en{Basically, we need more time.} \tabularnewline
\en{literally} & « littéralement » & Insistance (très familier, à doser) : \en{I literally waited two hours.} \tabularnewline
\en{to be honest} & « honnêtement » & Franchise : \en{To be honest, I didn't like it.} \tabularnewline
\en{I guess / I suppose} & « je suppose » & Hésitation polie : \en{I guess you're right.} \tabularnewline
\en{sort of} & « en quelque sorte » & Approximation. \tabularnewline
\en{like} & « genre » (très familier) & \en{It was, like, ten metres high.} --- à éviter en contexte formel. \tabularnewline
\en{anyway} & « bref, enfin » & Reprendre le fil : \en{Anyway, let's move on.} \tabularnewline
\en{so} & « donc, alors » & Enchaînement : \en{So, what happened next?} \tabularnewline
\en{right} & « d'accord, bon » & Validation : \en{Right, I understand now.} \tabularnewline
\en{OK} & « d'accord » & Transition. \tabularnewline
\en{er / erm} & « euh » & Le vrai « euh » anglais --- bref et rare. \tabularnewline
\en{let's see} & « voyons » & Réflexion. \tabularnewline
\en{you see} & « tu comprends » & Explication. \tabularnewline
\en{as I was saying} & « comme je disais » & Reprendre après interruption. \tabularnewline
\en{by the way} & « au fait » & Digression. \tabularnewline
\en{that said} & « cela dit » & Nuance. \tabularnewline
\en{on the other hand} & « d'un autre côté » & Contraste. \tabularnewline
\en{in fact} & « en fait » & Précision (attention : pas \en{actually} partout). \tabularnewline
\en{for example / for instance} & « par exemple » & Illustration. \tabularnewline
\en{at the end of the day} & « au bout du compte » & Conclusion. \tabularnewline
\end{tabularx}
\end{tableau}

\begin{piege}[Trois \en{fillers} maximum]
Trop de \en{you know} et de \en{like} donne l'impression inverse de la fluidité : un discours
hésitant. Choisissez \textbf{deux ou trois} marqueurs qui vous plaisent (\en{well},
\en{I mean}, \en{actually}, \en{to be honest}) et utilisez-les proprement. Le silence bref
vaut mieux qu'un \en{like} répété dix fois.
\end{piege}

\section{Les sept drills de fluidité}

\subsection{Drill 1 --- Les 60 secondes}
Choisissez un mot (\en{water}, \en{family}, \en{work}, \en{rain}) et parlez \textbf{60
secondes sans vous arrêter} sur ce mot : ce que c'est, ce que vous aimez, un souvenir,
une opinion. Chronométrez. Trois fois par jour. Au bout d'un mois, vous tiendrez
90 secondes sans blanc.

\subsection{Drill 2 --- La chaîne de mots}
Partez d'un mot, associez librement, et justifiez chaque enchaînement en une phrase :
\en{rain \ra umbrella \ra my grandmother \ra her house \ra the village\ldots}
Cinq minutes, à voix haute. Ce drill casse la peur de la phrase suivante.

\subsection{Drill 3 --- Le récit en trois temps}
Racontez la même anecdote \textbf{trois fois} : la première en 3 minutes, la deuxième en
2 minutes, la troisième en 1 minute. Vous apprenez à \textbf{condenser} et à éliminer le
superflu --- c'est l'exercice préféré des interprètes.

\subsection{Drill 4 --- Le double rôle}
Jouez seul un dialogue de ce livre, en faisant les deux voix, debout, avec gestes.
Changez une variable à chaque tour (vous êtes pressé, vous êtes poli, vous êtes agacé).

\subsection{Drill 5 --- La description d'image}
Prenez une photo (téléphone, journal, Internet) et décrivez-la pendant 2 minutes :
\en{In the foreground\ldots\ In the background\ldots\ It looks like\ldots\ They might be\ldots\
Maybe he's feeling\ldots} --- l'usage de \en{might} et \en{looks like} fait passer votre
anglais au niveau B2.

\subsection{Drill 6 --- La traduction différée}
Écoutez une phrase anglaise (podcast), mettez pause, \textbf{résumez-la en anglais} avec vos
propres mots. Jamais de traduction française : c'est ce qui installe la pensée directe.

\subsection{Drill 7 --- Le miroir à questions}
Préparez 30 questions sur votre vie (voir liste ci-dessous) et répondez-y à voix haute,
chacune en 45 secondes, sans préparation. Enregistrez. Répétez deux fois par semaine.

\begin{tableau}[30 questions sur vous-même --- le kit de l'entretien oral]
\footnotesize
\begin{multicols}{2}
\begin{anglais}
\begin{enumerate}\setlength{\itemsep}{1pt}
\item Can you tell me about yourself?
\item Where did you grow up?
\item What does your name mean?
\item Describe your family.
\item What do you do for a living?
\item Why did you choose that job?
\item What's a typical day like for you?
\item What did you do last weekend?
\item What are you planning to do next year?
\item Why are you learning English?
\item How do you learn best?
\item What was your favourite subject at school?
\item Describe your best friend.
\item What do you do in your free time?
\item What kind of music do you like?
\item Do you play any sport?
\item What's your favourite food, and how is it made?
\item Describe your city or village.
\item What's the best thing about your country?
\item What would you change about it?
\item Tell me about a memorable trip.
\item What was the happiest moment of your life?
\item Describe a difficulty you overcame.
\item Who do you admire, and why?
\item What are your three strengths?
\item What are you working on improving?
\item Where do you see yourself in five years?
\item What's your opinion on social media?
\item If you could learn anything, what would it be?
\item What advice would you give to a beginner learning English?
\end{enumerate}
\end{anglais}
\end{multicols}
\end{tableau}

\section{Fluidité contre précision : l'équilibre}
\begin{tableau}[Quand privilégier quoi]
\begin{tabularx}{\linewidth}{@{}p{4.6cm}p{4.6cm}X@{}}
\entete{Objectif} & \entete{Quand} & \entete{Comment} \tabularnewline
\hline
\textbf{Fluidité} & Conversation réelle, appels, discussions de groupe & Parlez sans vous arrêter, ne corrigez pas en chemin, acceptez les fautes, restez dans le rythme. \tabularnewline
\textbf{Précision} & Exercices écrits, préparation d'un exposé, enregistrement de travail & Corrigez chaque phrase, vérifiez les temps, les articles, la prononciation. \tabularnewline
\textbf{Règle 80/20} & Toute l'année & 80\,\% du temps de parole en mode fluide, 20\,\% en mode précis. L'inverse est l'erreur classique des francophones scolarisés. \tabularnewline
\end{tabularx}
\end{tableau}

\begin{astuce}[Ne vous corrigez jamais en pleine phrase]
S'interrompre pour corriger un \en{-s} oublié détruit le rythme, la confiance et le plaisir
de l'interlocuteur. Terminez votre idée, puis --- si la faute compte vraiment --- reformulez
la phrase entière correctement, une fois. Mieux : notez-la et corrigez-la le soir, hors
conversation.
\end{astuce}

\section{Le débit, les pauses, la voix}
\begin{itemize}
  \item \textbf{Débit :} visez 120 à 140 mots par minute en conversation. Trop lent =
  incompréhension chez l'autre ; trop rapide = fautes et finales mangées.
  \item \textbf{Pauses :} une pause courte après chaque groupe de sens. La pause n'est pas
  un échec, c'est de la ponctuation orale.
  \item \textbf{Volume :} parlez légèrement plus fort qu'en français. L'anglais porte.
  \item \textbf{Regard :} regardez votre interlocuteur, pas le plafond --- le regard fuyant
  fait croire que vous ne comprenez pas.
  \item \textbf{Respiration :} respirez par le ventre avant de commencer une phrase longue.
  Une phrase dite sans souffle se termine en murmure --- et l'information importante est perdue.
\end{itemize}

\section{Le plan anti-blocage en 6 réflexes}
\begin{enumerate}
  \item \textbf{Le mot manque ?} Contournez (\en{it's a thing you use to\ldots}).
  \item \textbf{La grammaire bloque ?} Simplifiez : sujet + verbe + complément, au présent.
  Une phrase simple et juste vaut mieux qu'une phrase complexe et fausse.
  \item \textbf{Vous n'avez pas compris ?} \en{Sorry, could you say that again more slowly?}
  --- c'est un réflexe à automatiser, pas une honte.
  \item \textbf{Vous paniquez ?} Buvez une gorgée d'eau, respirez, dites
  \en{Let me think for a second.} Le silence de trois secondes paraît long pour vous,
  court pour l'autre.
  \item \textbf{La conversation dérape vers un sujet que vous ne maîtrisez pas ?}
  Ramenez-la : \en{That reminds me of something I know well\ldots}
  \item \textbf{Vous avez fait une faute grave ?} Corrigez en souriant :
  \en{Sorry --- I mean\ldots} et continuez. Personne ne vous en tiendra rigueur.
\end{enumerate}

\begin{mission}[Chapitre 15 --- le défi des 21 jours de parole]
Pendant 21 jours, faites chaque jour, sans exception :\\
\textbf{1.} 3 minutes de description d'image (drill 5).\\
\textbf{2.} 3 minutes de « 60 secondes » sur trois mots différents (drill 1).\\
\textbf{3.} 2 minutes de réponse à trois questions de la liste des 30 (drill 7).\\
\textbf{4.} Enregistrez-vous deux fois dans la semaine et notez trois progrès.\\
Total : 8 minutes par jour. C'est le plus petit investissement du livre pour le plus grand
bénéfice : après 21 jours, les mots arriveront \textbf{avant} que vous ayez fini de les
chercher.
\end{mission}
